%% USPSC-Cap2-FundamentacaoTeorica.tex
% ----------------------------------------------------------
% Fundamentação Teórica
% ----------------------------------------------------------
\chapter{Fundamentação Teórica}
\label{cap:fundamentacao}

\section{Considerações Iniciais}
\label{sec:fund_consideracoes}

Este capítulo estabelece o arcabouço conceitual que sustenta a avaliação relatada
nesta monografia. A organização do texto segue a ordem em que os conceitos são
mobilizados no restante do trabalho. A Seção~\ref{sec:usabilidade_ux} delimita os
conceitos de usabilidade e de experiência do usuário a partir das normas
internacionais que os definem, e discute por que o contexto de uso móvel e ao ar
livre impõe restrições que não se manifestam em avaliações conduzidas em
laboratório. A Seção~\ref{sec:mobile} caracteriza o desenvolvimento móvel
multiplataforma e as implicações da escolha tecnológica sobre a qualidade da
interação. A Seção~\ref{sec:metodos} apresenta os métodos de avaliação
empregados --- avaliação heurística, protocolo \textit{Think-Aloud}, entrevista
semiestruturada e o \textit{Technology Acceptance Model} ---, bem como os
instrumentos considerados e deliberadamente não adotados. A
Seção~\ref{sec:personas_teoria} trata da modelagem de perfis de usuário por meio
de personas. A Seção~\ref{sec:gamificacao} discute o problema teórico específico
deste trabalho: a ambiguidade da origem do atrito em sistemas que incorporam
desafio deliberado. A Seção~\ref{sec:justificativa_mista} justifica a articulação
dos métodos em um desenho de métodos mistos.

Uma observação preliminar orienta a leitura do capítulo. A avaliação de um
artefato interativo não é procedimento neutro: cada método incorpora pressupostos
sobre a natureza do que se pretende medir e produz evidências de um tipo
particular. A escolha dos métodos aqui apresentados, portanto, não decorre de sua
popularidade na literatura, mas de sua adequação ao problema de pesquisa
enunciado na Seção~\ref{sec:problema} --- distinguir, em campo, as dificuldades
originadas na interface daquelas inerentes ao desafio proposto pela atividade.

\section{Usabilidade e Experiência do Usuário}
\label{sec:usabilidade_ux}

\subsection{Usabilidade segundo a ISO 9241-11}

A norma \citeonline{iso9241part11} define usabilidade como a medida em que um
sistema, produto ou serviço pode ser utilizado por usuários específicos para
alcançar objetivos especificados com \textbf{eficácia}, \textbf{eficiência} e
\textbf{satisfação} em um \textbf{contexto de uso} especificado. Cada termo da
definição carrega consequências metodológicas que convém explicitar.

\begin{alineas}
  \item \textbf{Eficácia} refere-se à acurácia e à completude com que os usuários
  atingem os objetivos. Operacionalmente, traduz-se em taxas de conclusão de
  tarefa e em contagens de erro.

  \item \textbf{Eficiência} relaciona os recursos consumidos --- tempo, esforço
  físico, esforço cognitivo --- aos resultados obtidos. É comumente aferida por
  tempo de tarefa e número de ações.

  \item \textbf{Satisfação} designa as respostas físicas, cognitivas e emocionais
  do usuário decorrentes do uso, e é tipicamente medida por instrumentos
  autorrelatados.

  \item \textbf{Contexto de uso} compreende os usuários, as tarefas, os
  equipamentos e o ambiente físico e social em que a interação ocorre.
\end{alineas}

A inclusão do contexto de uso na própria definição é o ponto de maior relevância
para esta monografia. A usabilidade não é atributo intrínseco do artefato, mas
propriedade relacional entre artefato, usuário, tarefa e ambiente. Segue-se que
uma avaliação conduzida em condições artificialmente favoráveis --- em ambiente
interno, com o usuário sentado e em repouso --- não caracteriza a usabilidade do
sistema no contexto em que ele será efetivamente utilizado. Esse raciocínio
fundamenta a decisão, detalhada no Capítulo~\ref{cap:metodologia}, de conduzir a
avaliação empírica em campo, no próprio parque, com padronização e registro das
condições ambientais em vez de sua eliminação.

\subsection{Experiência do Usuário segundo a ISO 9241-210}

A norma \citeonline{iso9241part210} amplia o escopo da avaliação ao definir
experiência do usuário como as percepções e respostas de uma pessoa resultantes
do uso ou da antecipação do uso de um sistema, produto ou serviço. A definição
abrange emoções, crenças, preferências, percepções, respostas físicas e
psicológicas, comportamentos e realizações que ocorrem antes, durante e depois do
uso.

A distinção entre usabilidade e experiência do usuário é frequentemente tratada
como questão de amplitude, mas é mais produtivo compreendê-la como diferença de
objeto. A usabilidade pergunta se o usuário consegue realizar a tarefa; a
experiência do usuário pergunta o que a realização da tarefa significou para ele.
No caso do \textit{Trilha das Árvores}, essa distinção é constitutiva do problema:
não basta verificar que o visitante conseguiu chegar ao \textit{checkpoint} e
escanear o QR Code; interessa saber se a mediação tecnológica enriqueceu ou
empobreceu sua relação com o parque.

\subsection{Qualidade em uso e qualidade de produto}

A norma \citeonline{isoiec25010} organiza a qualidade de sistemas de software em
dois modelos complementares. O modelo de \textbf{qualidade em uso} descreve os
efeitos do uso do sistema em um contexto específico, com características como
eficácia, eficiência, satisfação, ausência de risco e cobertura de contexto. O
modelo de \textbf{qualidade de produto} descreve propriedades estáticas e
dinâmicas do software, entre as quais a adequação funcional, a eficiência de
desempenho, a confiabilidade, a segurança e a usabilidade --- esta última
decomposta em subcaracterísticas como reconhecibilidade da adequação,
aprendibilidade, operabilidade, proteção contra erro do usuário, estética da
interface e acessibilidade.

A pertinência dessa norma para este trabalho está na articulação que ela permite
entre as duas camadas. A hipótese enunciada na Seção~\ref{sec:problema} sustenta
que os atritos percebidos pelo usuário --- fenômenos de qualidade em uso ---
originam-se, em parte, em decisões situadas na camada de qualidade de produto,
como a estratégia de persistência de imagens ou a política de tratamento de erros
de rede. A matriz de rastreabilidade proposta como contribuição do trabalho é,
sob essa ótica, um instrumento que vincula as duas camadas do modelo.

\subsection{Especificidades do contexto móvel e de uso em campo}
\label{sec:contexto_movel}

Aplicações móveis utilizadas em deslocamento operam sob restrições que sistemas
de mesa não enfrentam e que a literatura de Interação Humano-Computador tem
tratado sob a designação de limitações situacionais. Entre as mais relevantes
para este trabalho, destacam-se:

\begin{alineas}
  \item \textbf{Atenção dividida.} O usuário não dedica atenção exclusiva à tela.
  Ele caminha, observa o ambiente, desvia de obstáculos e mantém orientação
  espacial. A carga cognitiva disponível para a interface é, portanto, residual.

  \item \textbf{Condições de luminosidade não controladas.} A incidência direta
  de luz solar reduz drasticamente o contraste efetivo da tela, o que penaliza
  interfaces que dependem de diferenças cromáticas sutis ou de tipografia de
  pequeno corpo.

  \item \textbf{Interação motora degradada.} A precisão do toque diminui com o
  movimento do corpo, o que amplia o custo de alvos pequenos ou próximos entre
  si.

  \item \textbf{Conectividade e posicionamento intermitentes.} Áreas arborizadas
  atenuam o sinal de satélites de posicionamento e de redes móveis, o que produz
  latência e imprecisão que o usuário percebe como falha do aplicativo.

  \item \textbf{Restrição energética.} O uso contínuo de tela, GPS, câmera e
  sensores inerciais consome bateria em ritmo elevado, e a ansiedade quanto à
  autonomia do dispositivo integra a experiência.
\end{alineas}

Essas restrições não constituem meros incômodos a serem contornados no
planejamento do estudo. Elas são parte do contexto de uso previsto pela
\citeonline{iso9241part11} e, portanto, parte daquilo que se avalia. Sua
supressão descaracterizaria o objeto.

\section{Desenvolvimento Móvel Multiplataforma}
\label{sec:mobile}

\subsection{Estratégias de desenvolvimento}

A computação móvel consolidou-se como principal meio de acesso à informação, e o
mercado de sistemas operacionais móveis organizou-se em torno de duas plataformas
dominantes, Android e iOS. Essa configuração impõe ao desenvolvimento o problema
da fragmentação: historicamente, atender ambas as plataformas exigia manter duas
bases de código independentes, escritas em Kotlin ou Java e em Swift ou
Objective-C, respectivamente, com duplicação de esforço e ampliação da superfície
de inconsistências entre versões.

Em resposta a esse problema, consolidaram-se três estratégias, sintetizadas no
Quadro~\ref{quad:estrategias}.

\begin{quadro}[htb]
\centering
\caption{Estratégias de desenvolvimento de aplicações móveis}
\label{quad:estrategias}
\begin{tabular}{p{2.6cm} p{4.2cm} p{4.6cm}}
\toprule
\textbf{Estratégia} & \textbf{Característica} & \textbf{Implicação para a UX} \\
\midrule
Nativa & Código específico por plataforma, com acesso direto às APIs do sistema &
Desempenho e integração de sensores máximos; custo de desenvolvimento e
manutenção mais elevado \\
\addlinespace
Híbrida (\textit{web}) & Conteúdo web encapsulado em contêiner nativo &
Menor custo; latência superior e acesso limitado a sensores, o que compromete
interações dependentes de hardware \\
\addlinespace
Multiplataforma & Base de código única traduzida para componentes nativos reais &
Reaproveitamento de código com desempenho próximo ao nativo e acesso a sensores
por bibliotecas padronizadas \\
\bottomrule
\end{tabular}
\\[4pt]
\begin{flushleft}
Fonte: Elaborado pelo autor.
\end{flushleft}
\end{quadro}

\subsection{React Native e a plataforma Expo}

O \textit{framework} React Native, mantido pela Meta, situa-se na abordagem
multiplataforma. Ele permite descrever a interface de forma declarativa em
JavaScript, delegando a renderização a componentes nativos reais de cada
plataforma, e não a uma \textit{webview}. Um mesmo componente declarado no código
é, assim, traduzido para os elementos visuais nativos do Android e do iOS,
preservando as convenções de interação de cada sistema.

Essa propriedade tem consequência direta sobre a usabilidade. Ao herdar os
componentes nativos, a aplicação herda também os padrões de interação que o
usuário já domina em seu próprio dispositivo, o que favorece a heurística de
consistência e padrões discutida na Seção~\ref{sec:heuristica}. Sobre o React
Native, o projeto emprega a plataforma Expo, que fornece um fluxo de trabalho
gerenciado, abstrai parte substancial da configuração nativa e disponibiliza
módulos padronizados de acesso a hardware.

\subsection{Implicações da escolha tecnológica para a experiência}

A pertinência da abordagem multiplataforma para o \textit{Trilha das Árvores} não
é apenas econômica. A proposta de valor do aplicativo depende de integração de
baixa latência com três subsistemas de hardware: o receptor de posicionamento
global, que sustenta o mapa em tempo real e o cálculo de distância; o
magnetômetro, que alimenta a bússola; e a câmera, que realiza a leitura dos QR
Codes nos \textit{checkpoints}. Uma abordagem híbrida introduziria, em cada um
desses caminhos, uma camada adicional de tradução entre o contêiner web e as APIs
nativas, com custo em latência perceptível. Em uma atividade realizada em
movimento, latência perceptível converte-se em hesitação, e hesitação converte-se
em deslocamento de atenção do ambiente para a tela --- exatamente o efeito que o
projeto pretende evitar.

\section{Métodos de Avaliação de Usabilidade e Experiência do Usuário}
\label{sec:metodos}

\subsection{Métodos de inspeção e métodos empíricos}

A literatura de Interação Humano-Computador organiza os métodos de avaliação em
dois grandes grupos \cite{barbosasilva2010}. Os \textbf{métodos de inspeção},
também chamados analíticos, dispensam a participação de usuários: um ou mais
especialistas examinam a interface à luz de princípios, modelos ou padrões
consolidados. Os \textbf{métodos empíricos}, ou de investigação com usuários,
coletam dados da interação real ou das opiniões de usuários efetivos ou
potenciais.

Os dois grupos produzem evidências de naturezas distintas e não substituíveis. A
inspeção identifica violações de princípios com baixo custo e em qualquer estágio
do desenvolvimento, mas não informa se as violações identificadas efetivamente
prejudicam usuários reais, nem revela problemas que decorram de descompasso entre
o modelo mental do projetista e o do usuário. A investigação empírica revela
esses descompassos, mas é onerosa, exige um artefato em estado utilizável e, se
aplicada a uma interface ainda repleta de defeitos elementares, desperdiça a
atenção dos participantes em problemas que a inspeção teria detectado sem custo.

Dessa complementaridade decorre o encadeamento adotado neste trabalho: a
inspeção heurística precede e sanea, e o estudo empírico verifica e explica.

\subsection{Avaliação Heurística}
\label{sec:heuristica}

A Avaliação Heurística, proposta por \citeonline{nielsenmolich1990} e consolidada
por \citeonline{nielsen1994heuristic}, é o método de inspeção mais difundido na
engenharia de usabilidade. O procedimento consiste em um pequeno conjunto de
avaliadores examinando a interface individualmente, percorrendo seus elementos e
registrando toda violação de um conjunto de princípios gerais de usabilidade ---
as heurísticas. Os relatos individuais são posteriormente consolidados em uma
lista única de problemas.

O conjunto de referência é o das dez heurísticas de
\citeonline{nielsen1994heuristic}, reproduzidas no Quadro~\ref{quad:heuristicas}.

\begin{quadro}[htb]
\centering
\caption{As dez heurísticas de usabilidade}
\label{quad:heuristicas}
\begin{tabular}{c p{4.0cm} p{6.4cm}}
\toprule
\textbf{\#} & \textbf{Heurística} & \textbf{Descrição} \\
\midrule
H1 & Visibilidade do status do sistema & O sistema mantém o usuário informado
sobre o que ocorre, por meio de retorno apropriado em tempo razoável \\
\addlinespace
H2 & Correspondência entre o sistema e o mundo real & O sistema emprega a
linguagem do usuário, com palavras e conceitos que lhe são familiares, em ordem
natural e lógica \\
\addlinespace
H3 & Controle e liberdade do usuário & O sistema oferece saídas claramente
sinalizadas para estados indesejados, com desfazer e refazer \\
\addlinespace
H4 & Consistência e padrões & Termos, situações e ações equivalentes são
representados de forma equivalente, respeitando as convenções da plataforma \\
\addlinespace
H5 & Prevenção de erros & O projeto elimina condições propensas a erro ou
apresenta confirmação antes que o usuário se comprometa com a ação \\
\addlinespace
H6 & Reconhecimento em vez de memorização & Objetos, ações e opções permanecem
visíveis, dispensando que o usuário retenha informação entre telas \\
\addlinespace
H7 & Flexibilidade e eficiência de uso & Aceleradores, invisíveis ao usuário
novato, permitem que o usuário experiente atue com mais rapidez \\
\addlinespace
H8 & Estética e projeto minimalista & As interfaces não contêm informação
irrelevante ou raramente necessária, que compete com a informação pertinente \\
\addlinespace
H9 & Ajuda ao usuário no reconhecimento, diagnóstico e recuperação de erros &
As mensagens de erro são expressas em linguagem simples, indicam precisamente o
problema e sugerem uma solução \\
\addlinespace
H10 & Ajuda e documentação & Quando necessária, a documentação é fácil de
localizar, focada na tarefa do usuário e concisa \\
\bottomrule
\end{tabular}
\\[4pt]
\begin{flushleft}
Fonte: Adaptado de \citeonline{nielsen1994heuristic}.
\end{flushleft}
\end{quadro}

Cada problema identificado recebe uma classificação de severidade, que combina a
frequência com que o problema ocorre, o impacto que produz sobre o usuário e sua
persistência ao longo do uso repetido. A escala empregada é a de
\citeonline{nielsen1994heuristic}, apresentada no
Quadro~\ref{quad:severidade}, e sua função é ordenar as correções segundo
prioridade, em um cenário de recursos de desenvolvimento limitados.

\begin{quadro}[htb]
\centering
\caption{Escala de severidade de problemas de usabilidade}
\label{quad:severidade}
\begin{tabular}{c p{9.6cm}}
\toprule
\textbf{Grau} & \textbf{Interpretação} \\
\midrule
0 & Não se trata de problema de usabilidade \\
1 & Problema cosmético; corrigir apenas se houver tempo disponível \\
2 & Problema menor; correção de baixa prioridade \\
3 & Problema maior; correção de alta prioridade \\
4 & Catastrófico; correção obrigatória antes da liberação \\
\bottomrule
\end{tabular}
\\[4pt]
\begin{flushleft}
Fonte: Adaptado de \citeonline{nielsen1994heuristic}.
\end{flushleft}
\end{quadro}

Quanto ao número de avaliadores, \citeonline{nielsenlandauer1993} propuseram um
modelo segundo o qual a proporção de problemas encontrados cresce de forma
assintótica com o número de inspetores, com ganho marginal decrescente a partir
de aproximadamente cinco avaliadores. A recomendação prática daí derivada ---
empregar de três a cinco avaliadores --- é a que orienta o dimensionamento
descrito no Capítulo~\ref{cap:metodologia}.

Reconhecem-se as limitações do método. A avaliação heurística tende a produzir
falsos positivos, isto é, apontamentos que não se confirmam como problemas para
usuários reais; é sensível à experiência do avaliador, com especialistas duplos
--- conhecedores tanto de usabilidade quanto do domínio --- identificando
substancialmente mais problemas; e é pouco eficaz na detecção de problemas que
dependam de conhecimento tácito do domínio ou de expectativas específicas do
público-alvo. Essas limitações são precisamente as que a etapa empírica se
destina a compensar.

\subsection{Percurso Cognitivo}

O Percurso Cognitivo, proposto por \citeonline{wharton1994}, constitui método de
inspeção alternativo, centrado na aprendibilidade. O avaliador percorre uma
tarefa passo a passo e, a cada ação, formula quatro perguntas relativas à
possibilidade de o usuário formar o objetivo correto, localizar o controle
adequado, reconhecê-lo como o controle correto e interpretar o retorno obtido.

O método foi considerado no planejamento deste trabalho e não foi adotado como
técnica principal, por três razões. Primeiro, seu foco é o uso exploratório por
usuário novato, dimensão que o estudo empírico com participantes em primeiro
contato já cobre com evidência direta em vez de simulada. Segundo, sua aplicação
rigorosa é demorada e produz recorte estreito, restrito às tarefas percorridas.
Terceiro, o problema central desta pesquisa --- a atribuição da origem do atrito
--- exige acesso ao processo cognitivo do usuário real, e não à simulação que o
avaliador faz dele.

\subsection{Protocolo Think-Aloud}
\label{sec:thinkaloud}

\subsubsection{Fundamentação psicológica}

O protocolo \textit{Think-Aloud} consiste em solicitar ao participante que
verbalize continuamente seus pensamentos enquanto executa uma tarefa. Sua
fundamentação teórica remonta a \citeonline{ericssonsimon1980} e ao tratamento
sistemático posteriormente consolidado em \citeonline{ericssonsimon1993}, que
distinguem três níveis de verbalização:

\begin{alineas}
  \item \textbf{Nível 1 --- vocalização direta.} O participante enuncia
  informação já codificada verbalmente na memória de trabalho. Não há
  processamento adicional e não há interferência sobre o desempenho da tarefa.

  \item \textbf{Nível 2 --- codificação verbal.} O participante traduz em
  palavras informação armazenada em outro formato, como imagens ou relações
  espaciais. Há custo cognitivo adicional, mas o conteúdo do pensamento permanece
  inalterado.

  \item \textbf{Nível 3 --- explicação e justificativa.} O participante é
  solicitado a explicar, justificar ou racionalizar suas ações. Esse nível exige
  processos cognitivos que não fariam parte da execução da tarefa e,
  consequentemente, \textbf{altera} tanto o processo quanto o desempenho.
\end{alineas}

A consequência metodológica dessa distinção é operacional e rígida: as instruções
dadas ao participante e as intervenções do moderador devem restringir-se a
solicitar verbalização de níveis 1 e 2. Perguntas como ``por que você fez isso?''
induzem verbalização de nível 3 e comprometem a validade dos dados coletados. O
roteiro de moderação apresentado no Capítulo~\ref{cap:metodologia} incorpora essa
restrição de forma explícita, com uma lista fechada de sondas permitidas.

\subsubsection{Verbalização concomitante e retrospectiva}

Distinguem-se duas modalidades de aplicação. Na modalidade
\textbf{concomitante}, a verbalização ocorre simultaneamente à execução da
tarefa; seu conteúdo reflete o estado corrente da memória de trabalho, sem os
vieses de racionalização e de esquecimento. Na modalidade
\textbf{retrospectiva}, o participante relata seus pensamentos após a conclusão
da tarefa, frequentemente auxiliado por um estímulo de recuperação, como a
reprodução do vídeo da sessão.

\citeonline{vandenhaak2003} compararam experimentalmente as duas modalidades e
observaram um compromisso entre elas: a modalidade concomitante produz maior
volume de dados sobre o processo, ao custo de interferir na execução da tarefa;
a retrospectiva preserva o desempenho e permite comentários mais elaborados, mas
está sujeita a lacunas de memória e à reconstrução \textit{a posteriori}. Essa
constatação é o fundamento da adaptação descrita a seguir.

\subsubsection{Reatividade, moderação e limites do método}

A crítica de \citeonline{borenramey2000} é particularmente relevante para a
prática. Os autores observam que a aplicação corrente do método em avaliações de
usabilidade afasta-se substancialmente do protocolo estrito de Ericsson e Simon,
e propõem uma fundamentação alternativa, baseada na teoria da comunicação: a
sessão é um evento comunicativo entre dois participantes, no qual o moderador
sustenta a fala do usuário por meio de sinais mínimos de escuta, sem interferir no
conteúdo. Dessa proposição derivam recomendações concretas adotadas neste
trabalho: manter presença audível por meio de marcadores de reconhecimento,
intervir com fórmulas fixas e neutras, e jamais responder a perguntas do
participante durante a execução.

Reconhecem-se, por rigor, as limitações do protocolo. A verbalização pode alterar
levemente o desempenho, fenômeno designado reatividade; exige do participante
esforço adicional, agravado quando a tarefa envolve deslocamento físico; e demanda
do moderador disciplina considerável para não induzir ações. Essas limitações são
mitigadas por treinamento prévio de verbalização, por roteiro fechado de
intervenções e pela triangulação com os demais instrumentos.

\subsubsection{Adaptação do protocolo ao contexto de campo}

A aplicação do protocolo concomitante durante todo o percurso de uma trilha
apresenta um problema específico. A fala exige recurso respiratório e atencional
que compete com a caminhada, o que degrada simultaneamente a qualidade da
verbalização e o desempenho na atividade --- e, em um aplicativo cuja proposta
inclui componente de esforço físico, essa degradação confunde-se com o próprio
objeto de estudo.

Adota-se, por essa razão, um \textbf{protocolo híbrido segmentado}, que alterna
as duas modalidades conforme o segmento da atividade:

\begin{alineas}
  \item nos segmentos \textbf{estacionários}, em que o participante manipula a
  interface parado, aplica-se verbalização concomitante;

  \item nos segmentos de \textbf{deslocamento} entre \textit{checkpoints},
  observa-se o comportamento sem exigir fala, e aplica-se recuperação
  retrospectiva imediata na chegada, enquanto o traço em memória permanece
  acessível;

  \item ao final da sessão, aplica-se recuperação retrospectiva assistida por
  estímulo, com reprodução da gravação de tela dos episódios de hesitação
  previamente registrados.
\end{alineas}

Essa articulação preserva a validade ecológica da observação em campo e mantém a
qualidade dos dados verbais, apoiando-se no compromisso entre modalidades
identificado por \citeonline{vandenhaak2003}.

% ============================================================================
% >>> [PREENCHER --- SONDAS E PERGUNTAS DO THINK-ALOUD] <<<
%
% Espaço reservado, conforme planejado, para a descrição das perguntas e
% sondas utilizadas durante o protocolo.
%
% Sugestão de conteúdo a redigir aqui:
%   (a) a lista fechada de sondas permitidas durante a tarefa, com a
%       justificativa de nível de verbalização de cada uma;
%   (b) a lista de formulações proibidas e o motivo (indução de nível 3);
%   (c) a fórmula de devolução usada quando o participante faz perguntas;
%   (d) o critério temporal de intervenção no silêncio;
%   (e) o protocolo de ajuda escalonada e o critério de abandono de tarefa.
%
% Observação: o roteiro operacional completo entra como APÊNDICE. Aqui cabe
% apenas a fundamentação de cada escolha, remetendo ao apêndice.
% ============================================================================

\subsection{Entrevista semiestruturada pós-tarefa}
\label{sec:entrevista_teoria}

A verbalização concomitante, por sua própria natureza, não captura aquilo que o
participante não processa conscientemente durante a tarefa. Elementos de
interface que ele simplesmente evita --- um ícone cujo significado não reconhece,
um controle que não identifica como acionável --- não geram verbalização, pois não
chegam a integrar seu fluxo de pensamento. O silêncio, nesses casos, é ausência
de dado e não ausência de problema.

A entrevista semiestruturada aplicada imediatamente após as tarefas destina-se a
recuperar essa camada. Sua estrutura adota o formato de funil: inicia-se com
perguntas abertas e não dirigidas, que capturam o que é saliente para o
participante sem sugestão do entrevistador, e progride para sondagens dirigidas
às hipóteses do estudo, incluindo verificação explícita de reconhecimento de
ícones e de interpretação de indicadores. A ordenação não é acessória: perguntas
dirigidas formuladas antes das abertas contaminam as respostas subsequentes.

\subsection{Modelo de Aceitação de Tecnologia (TAM)}
\label{sec:tam_teoria}

O \textit{Technology Acceptance Model}, proposto por \citeonline{davis1989}, é uma
adaptação da Teoria da Ação Racional ao domínio da aceitação de sistemas
computacionais. O modelo sustenta que o uso efetivo de uma tecnologia é
determinado pela intenção comportamental de uso, a qual, por sua vez, é
determinada por dois construtos:

\begin{alineas}
  \item \textbf{Utilidade Percebida} (\textit{Perceived Usefulness} --- PU): o
  grau em que a pessoa acredita que o uso de determinado sistema melhoraria seu
  desempenho;

  \item \textbf{Facilidade de Uso Percebida} (\textit{Perceived Ease of Use} ---
  PEOU): o grau em que a pessoa acredita que o uso de determinado sistema seria
  livre de esforço.
\end{alineas}

O modelo postula ainda que a PEOU exerce efeito sobre a PU: um sistema mais fácil
de operar tende a ser percebido como mais útil, uma vez que menor parcela do
esforço do usuário é consumida pela operação e maior parcela pode ser dirigida ao
objetivo. Extensões posteriores, entre as quais o TAM2 e o modelo integrado
apresentado por \citeonline{venkateshbala2008}, incorporaram determinantes
adicionais de ambos os construtos, como norma subjetiva, relevância para o
trabalho, autoeficácia e percepção de controle externo.

A pertinência do TAM para este trabalho apoia-se em três argumentos. Primeiro, o
objetivo do projeto transcende a operabilidade: interessa saber se o público
adotaria e continuaria a usar o aplicativo, o que faz da Intenção de Uso um
desfecho legítimo, ausente de instrumentos que meçam apenas usabilidade. Segundo,
os dois construtos centrais correspondem com precisão à tese defendida na
Seção~\ref{sec:motivacao}: a PEOU operacionaliza a exigência de que a interface
não desvie carga cognitiva do desafio ambiental, enquanto a PU verifica se a
mediação agrega valor efetivo à exploração. Terceiro, o modelo é parcimonioso e
produz escores comparáveis, adequados a um estudo de graduação e a rodadas
sucessivas de avaliação.

A operacionalização emprega escala psicométrica do tipo Likert
\cite{likert1932}. Adota-se amplitude de sete pontos, com ponto neutro central.
A escolha de sete níveis, em vez de três ou cinco, oferece maior granularidade e
sensibilidade a variações de intensidade e reduz a concentração de respostas em
categorias extremas por ausência de alternativa intermediária adequada. Os itens
efetivamente aplicados são apresentados na Seção~\ref{sec:instrumento_tam}.

\subsection{Instrumentos considerados e não adotados}
\label{sec:instrumentos_excluidos}

A decisão por um instrumento de aceitação, em detrimento de escalas de usabilidade
percebida ou de qualidade hedônica, exige justificativa explícita. Foram
considerados e não adotados os três instrumentos a seguir.

O \textbf{System Usability Scale} \cite{brooke1996} é escala de dez itens,
amplamente empregada por sua brevidade e por produzir escore único de
interpretação consolidada. Não foi adotado como instrumento principal porque mede
usabilidade percebida de forma agregada, sem discriminar utilidade de facilidade,
e não contempla intenção de uso --- dimensão que, como argumentado, integra o
objetivo do projeto. Registra-se que sua aplicação complementar em rodadas
futuras permanece viável e desejável.

O \textbf{User Experience Questionnaire} \cite{laugwitz2008} e o
\textbf{AttrakDiff} \cite{hassenzahl2003} avaliam, além das qualidades
pragmáticas, as qualidades hedônicas da interação --- estímulo, identificação,
originalidade, atratividade. Não foram adotados porque o caráter do software é
utilitário e mediador: em uma atividade de exploração física no mundo real, o
estímulo e o engajamento devem residir no ambiente e no desafio, e não na fruição
da interface. Reconhece-se que essa opção prioriza a dimensão pragmática da
experiência, e a investigação sistemática das qualidades hedônicas fica
registrada como perspectiva de trabalho futuro.

\section{Personas e Proto-personas}
\label{sec:personas_teoria}

Persona é a descrição de um arquétipo de usuário, construída como personagem
fictício dotado de nome, objetivos, comportamentos, contexto e frustrações,
destinada a orientar decisões de projeto e de avaliação. A técnica foi
popularizada por \citeonline{cooper2004} e sistematizada por
\citeonline{pruittgrudin2003}, que discutem tanto seu poder de comunicação quanto
os riscos de sua construção sem base empírica.

A persona clássica é derivada de pesquisa com usuários reais --- entrevistas,
observação etnográfica, dados de uso. Quando essa base não está disponível,
recorre-se à \textbf{proto-persona} \cite{gothelf2013}: construção que explicita
as premissas da equipe sobre o público-alvo, formulada como hipótese a ser
verificada e revisada à medida que dados empíricos se acumulem.

A distinção é adotada aqui deliberadamente. O presente trabalho não dispõe de
pesquisa prévia sobre o perfil de visitação do parque, e apresentar como personas
o que são premissas seria atribuir aos perfis uma autoridade empírica que eles não
possuem. Adotam-se, portanto, proto-personas, explicitamente identificadas como
tal, cuja função é dupla: orientar a definição dos critérios de recrutamento e o
desenho das tarefas, e constituir objeto de verificação, uma vez que os dados
coletados permitirão confirmá-las, corrigi-las ou refutá-las. Registra-se ainda o
uso de uma \textbf{antipersona}, isto é, a caracterização explícita de um perfil
que o produto não se destina a atender, cuja função metodológica é fundamentar os
critérios de exclusão do recrutamento.

\section{Gamificação e a Ambiguidade da Origem do Atrito}
\label{sec:gamificacao}

Entende-se por gamificação o uso de elementos de projeto característicos de jogos
em contextos que não são jogos \cite{deterding2011}. O \textit{Trilha das
Árvores} incorpora vários desses elementos: um percurso com \textit{checkpoints}
sequenciais, validação de chegada por leitura de código, cronometragem, medição de
distância e indicação de progresso.

Essa incorporação cria um problema teórico que a literatura de usabilidade não
resolve de forma imediata. Os métodos clássicos de avaliação pressupõem que toda
dificuldade encontrada pelo usuário constitui defeito a ser eliminado, pois a
eficiência é tratada como valor a maximizar. Em sistemas com componente lúdico,
porém, parte da dificuldade é constitutiva da proposta: localizar a próxima árvore
exige esforço de orientação espacial, e esse esforço \textbf{é} a atividade, não um
obstáculo a ela. Um instrumento que registre tal esforço como falha de eficiência
produzirá recomendação de projeto que, se implementada, destruiria a proposta de
valor do produto.

Formula-se, assim, a necessidade de um critério de discriminação. Este trabalho
adota a classificação de cada episódio de atrito segundo sua \textbf{origem}, em
quatro categorias mutuamente excludentes:

\begin{alineas}
  \item \textbf{Interface} --- o atrito decorre de rótulo ambíguo, controle não
  reconhecido, retorno ausente ou fluxo inconsistente. É defeito e demanda
  correção.

  \item \textbf{Desafio proposto} --- o atrito decorre da atividade que o
  aplicativo deliberadamente propõe: caminhar, orientar-se, localizar a placa. É
  característica e não demanda correção, embora possa demandar melhor
  comunicação.

  \item \textbf{Infraestrutura} --- o atrito decorre de latência de rede, perda
  de sinal de posicionamento ou indisponibilidade de serviço externo. É defeito de
  natureza não interacional, cuja correção situa-se na camada de engenharia.

  \item \textbf{Dívida técnica} --- o atrito decorre de funcionalidade incompleta
  ou de dado não integrado, conhecidos pela equipe. É defeito de completude.
\end{alineas}

Essa classificação, aplicada a cada episódio registrado, é o elemento que
operacionaliza a questão de pesquisa formulada na Seção~\ref{sec:problema}. Sua
aplicação exige acesso à atribuição causal do próprio participante, o que
constitui argumento adicional em favor do protocolo \textit{Think-Aloud} e da
entrevista pós-tarefa, e contra instrumentos que produzam apenas medidas
agregadas de desempenho.

\section{Triangulação e Justificativa da Abordagem Mista}
\label{sec:justificativa_mista}

A articulação dos três métodos não é aditiva, mas triangular
\cite{creswell2011}. Cada um responde a uma pergunta distinta e é vulnerável a um
tipo distinto de erro; sua convergência é o que confere credibilidade aos
achados.

\begin{alineas}
  \item A \textbf{avaliação heurística} atua \textit{a priori} e responde à
  pergunta ``a interface viola princípios consolidados?''. É vulnerável a falsos
  positivos.

  \item O \textbf{protocolo Think-Aloud}, complementado pela entrevista, atua
  durante e imediatamente após o uso, e responde à pergunta ``por que o usuário
  encontrou dificuldade, e de onde ela veio?''. É vulnerável à reatividade e à
  pequena dimensão da amostra.

  \item O \textbf{instrumento TAM} atua ao final e responde à pergunta ``quanto de
  utilidade, facilidade e intenção de uso o sistema produziu?''. É vulnerável a
  vieses de autorrelato e de desejabilidade social.
\end{alineas}

A convergência esperada entre as evidências é verificável e, portanto,
falsificável: um atrito recorrentemente verbalizado no \textit{Think-Aloud} deve
encontrar correspondência em escores mais baixos de Facilidade de Uso Percebida, e
um construto bem avaliado no TAM deve corresponder à ausência de atritos
correlatos na observação. A ausência dessa correspondência é, ela própria, achado
relevante, pois indicaria que o instrumento quantitativo não capturou o fenômeno
observado --- ou que a verbalização registrou incômodo que o participante não
considera relevante o bastante para afetar sua avaliação global.

É essa articulação que sustenta a contribuição pretendida e que se materializa na
matriz de rastreabilidade descrita no Capítulo~\ref{cap:metodologia}.

\section{Trabalhos Relacionados}
\label{sec:trabalhos_relacionados}

% ============================================================================
% >>> [PREENCHER --- REVISÃO DE TRABALHOS CORRELATOS] <<<
%
% Esta seção deve reunir de 4 a 8 trabalhos correlatos e terminar com um
% quadro comparativo posicionando esta monografia entre eles.
%
% Eixos de busca sugeridos:
%   - avaliação de usabilidade de aplicativos móveis de trilha, turismo ou
%     patrimônio natural;
%   - aplicações de mediação em parques, jardins botânicos e arboretos;
%   - avaliação de UX em contexto de campo / mobile HCI in the wild;
%   - gamificação aplicada à educação ambiental;
%   - aplicação do TAM a aplicativos móveis de turismo ou educação.
%
% Bases sugeridas: ACM DL, IEEE Xplore, Scopus, Google Scholar; e, para
% produção nacional, os anais do IHC (SBC) e do SBIE.
%
% Sugestão de quadro comparativo com colunas:
%   Trabalho | Domínio | Métodos de avaliação | Uso em campo? | n |
%   Rastreia atrito até decisão técnica?
%
% A última coluna é o que posiciona esta monografia como contribuição.
% ============================================================================

\section{Considerações Finais}
\label{sec:fund_finais}

Este capítulo delimitou os conceitos de usabilidade e de experiência do usuário a
partir das normas que os definem, destacando que o contexto de uso integra a
própria definição de usabilidade e que, por consequência, a avaliação de uma
aplicação de uso em campo deve preservar --- e não suprimir --- as condições
reais de operação. Caracterizou o desenvolvimento móvel multiplataforma e as
implicações da escolha tecnológica sobre a latência de acesso a sensores.
Apresentou os métodos de avaliação adotados, sua fundamentação e suas limitações,
bem como os instrumentos considerados e não adotados, com justificativa explícita.
Discutiu a modelagem de perfis por proto-personas e formulou o problema teórico
central deste trabalho: a ambiguidade da origem do atrito em sistemas que
incorporam desafio deliberado, para a qual se propôs uma classificação em quatro
categorias.

O capítulo seguinte descreve o ecossistema \textit{Trilha das Árvores}, sua
arquitetura e as decisões de engenharia cujo efeito sobre a experiência do usuário
constitui objeto da avaliação.
